\documentclass[../DoAn.tex]{subfiles}
\begin{document}

Chương 3 đã trình bày nền tảng lý thuyết và các công nghệ được lựa chọn cho từng tầng trong kiến trúc hệ thống, bao gồm phần cứng IoT, giao thức LoRa, công nghệ phần mềm phía máy chủ và giao diện, lý thuyết tính toán AQI, cùng hạ tầng triển khai. Trên cơ sở đó, Chương 4 trình bày chi tiết quá trình thiết kế, triển khai và đánh giá hệ thống. Nội dung chương bao gồm: (i) thiết kế kiến trúc trong phần \ref{section:4.1}, (ii) thiết kế chi tiết trong phần \ref{section:4.2}, (iii) xây dựng ứng dụng trong phần \ref{section:4.3}, (iv) kiểm thử trong phần \ref{section:4.4}, và (v) triển khai trong phần \ref{section:4.5}.

\section{Thiết kế kiến trúc}
\label{section:4.1}

\subsection{Lựa chọn kiến trúc phần mềm}
\label{subsection:4.1.1}

Hệ thống được thiết kế theo kiến trúc bốn tầng, bao gồm tầng cảm biến (Sensor Layer), tầng trung chuyển (Gateway Layer), tầng máy chủ (Server Layer) và tầng hiển thị (Presentation Layer). Kiến trúc bốn tầng được lựa chọn vì ba lý do chính. Thứ nhất, mỗi tầng có trách nhiệm rõ ràng và có thể phát triển độc lập: tầng cảm biến thu thập dữ liệu môi trường, tầng trung chuyển chuyển tiếp dữ liệu từ mạng \gls{LoRa} sang Internet, tầng máy chủ xử lý nghiệp vụ và lưu trữ, tầng hiển thị cung cấp giao diện cho người dùng. Thứ hai, kiến trúc này cho phép mở rộng theo chiều ngang tại từng tầng — ví dụ thêm Sensor Node mới mà không thay đổi phần mềm Backend. Thứ ba, độ phức tạp phù hợp với quy mô hệ thống: triển khai trên một máy chủ VPS duy nhất, không yêu cầu kiến trúc Microservice.

\textbf{Tầng cảm biến} gồm các Sensor Node, mỗi node tích hợp vi điều khiển ESP32, ba module cảm biến (PMS7003, CCS811, AHT10) và module \gls{LoRa} AS32-TTL-100. Firmware sử dụng FreeRTOS để quản lý đồng thời bốn tác vụ: đọc cảm biến, gửi dữ liệu \gls{LoRa}, hiển thị OLED và đo pin. Dữ liệu được đóng gói thành gói tin nhị phân 18 byte và phát qua \gls{LoRa} 433 MHz.

\textbf{Tầng trung chuyển} gồm các Gateway, mỗi Gateway tích hợp ESP32 với module \gls{LoRa} và kết nối WiFi. Gateway nhận gói tin \gls{LoRa}, lưu vào bộ đệm vòng (ring buffer), sau đó gửi lên tầng máy chủ dưới dạng \gls{HTTP} POST (\gls{JSON} batch). Ở tầng này, hệ thống áp dụng mô hình Hub-Spoke: nhiều Sensor Node (spoke) gửi dữ liệu đến một Gateway (hub) trung tâm, cho phép một Gateway phục vụ nhiều Sensor Node trong phạm vi phủ sóng mà không cần định tuyến phức tạp.

\textbf{Tầng máy chủ} bao gồm Backend Server (Node.js/Express) và cơ sở dữ liệu (PostgreSQL với TimescaleDB và PostGIS), được đóng gói trong Docker container. Backend Server được tổ chức theo mô hình phân lớp gồm ba lớp: lớp Controller tiếp nhận \gls{HTTP} request, lớp Service chứa logic nghiệp vụ, và lớp Data Access tương tác với cơ sở dữ liệu thông qua Prisma \gls{ORM}. Ngoài ra, tầng máy chủ còn có Nginx đóng vai trò reverse proxy và phục vụ tệp tĩnh, cùng các tác vụ nền (cron job) kiểm tra thiết bị offline và dọn dẹp dữ liệu cũ.

\textbf{Tầng hiển thị} là ứng dụng web \gls{SPA} xây dựng bằng React 19 và Vite 8, giao tiếp với tầng máy chủ qua \gls{REST} \gls{API} (\gls{HTTP}/HTTPS) và Socket.IO (WebSocket). Tầng này cung cấp hai nhóm giao diện: dashboard cho người dùng (bản đồ AQI, biểu đồ lịch sử, xếp hạng) và trang quản trị cho quản trị viên (quản lý thiết bị, cảnh báo, cấu hình, xuất dữ liệu).

\subsection{Thiết kế tổng quan}
\label{subsection:4.1.2}

%todo: Biểu đồ gói UML tổng quan (package diagram) — user tự thêm hình

Hệ thống được tổ chức thành bốn nhóm gói chính tương ứng với bốn tầng kiến trúc, mỗi nhóm chứa các gói con với trách nhiệm cụ thể.

\textbf{Nhóm gói Tầng cảm biến (sensor-node)} chứa các module: \texttt{drivers} (giao tiếp cảm biến PMS7003, CCS811, AHT10, màn hình OLED, module LoRa), \texttt{core} (quản lý cấu hình NVS, provisioning qua WiFi AP, vòng lặp chính), và \texttt{config} (định nghĩa chân GPIO, tham số LoRa). Hệ thống có ba biến thể firmware cho các loại board ESP32 khác nhau.

\textbf{Nhóm gói Tầng trung chuyển (gateway)} chứa các module: \texttt{drivers} (giao tiếp LoRa), \texttt{core} (bộ đệm vòng, gửi \gls{HTTP} batch lên tầng máy chủ), và \texttt{config} (cấu hình WiFi, địa chỉ Backend). Nhóm gói này hoạt động độc lập, không phụ thuộc vào nhóm gói tầng cảm biến.

\textbf{Nhóm gói Tầng máy chủ (backend)} gồm sáu gói: \texttt{routes} (định nghĩa endpoint API), \texttt{controllers} (tiếp nhận request, trả response), \texttt{services} (logic nghiệp vụ), \texttt{middlewares} (xác thực JWT, xử lý lỗi), \texttt{validations} (kiểm tra dữ liệu đầu vào), và \texttt{config} (kết nối Prisma Client). Luồng phụ thuộc tuân thủ nguyên tắc một chiều: routes $\rightarrow$ controllers $\rightarrow$ services $\rightarrow$ config (Prisma), đảm bảo gói tầng dưới không phụ thuộc gói tầng trên.

\textbf{Nhóm gói Tầng hiển thị (frontend)} gồm bảy gói: \texttt{pages} (các trang giao diện), \texttt{components} (thành phần tái sử dụng), \texttt{hooks} (custom React hooks), \texttt{services} (gọi API), \texttt{stores} (quản lý trạng thái Zustand), \texttt{layouts} (bố cục trang), và \texttt{i18n} (đa ngôn ngữ). Gói \texttt{pages} phụ thuộc \texttt{components} và \texttt{hooks}; gói \texttt{hooks} phụ thuộc \texttt{services}; gói \texttt{services} gọi \gls{REST} \gls{API} đến tầng máy chủ.

\subsection{Thiết kế chi tiết gói}
\label{subsection:4.1.3}

%todo: Biểu đồ quan hệ class/module cho backend services — user tự thêm hình

Phần này trình bày thiết kế chi tiết cho hai nhóm gói quan trọng: nhóm xử lý dữ liệu telemetry trên Backend và nhóm hiển thị dashboard trên Frontend.

\textbf{Nhóm xử lý dữ liệu telemetry} trên Backend bao gồm bốn module liên kết chặt chẽ. Module \texttt{telemetryController} nhận HTTP POST từ Gateway, xác thực dữ liệu đầu vào qua \texttt{telemetryValidation}, sau đó chuyển cho \texttt{telemetryService}. Module \texttt{telemetryService} thực hiện ba bước trong một transaction: (i) lọc bỏ gói tin trùng lặp (deduplication) khi nhiều Gateway cùng nhận một gói LoRa, (ii) cập nhật trạng thái Gateway và Sensor Node, (iii) lưu dữ liệu đo lường vào bảng \texttt{measurements}. Sau transaction, module gọi \texttt{alertService.checkThresholdsAndCreateAlerts} để kiểm tra ngưỡng cảnh báo — việc tách riêng kiểm tra cảnh báo ra ngoài transaction đảm bảo lỗi cảnh báo không ảnh hưởng đến luồng lưu trữ dữ liệu chính. Module \texttt{aqiService} cung cấp hàm \texttt{calculateAQI} được gọi bởi \texttt{stationService} khi trả dữ liệu dashboard cho Frontend.

\textbf{Nhóm hiển thị dashboard} trên Frontend bao gồm trang \texttt{Dashboard} sử dụng custom hook \texttt{useDashboard} để gọi API lấy dữ liệu trạm gần nhất, và hook \texttt{useSocket} để nhận cập nhật thời gian thực qua Socket.IO. Khi Backend phát sự kiện \texttt{new\_telemetry\_data}, hook \texttt{useSocket} cập nhật trạng thái React, kích hoạt Virtual DOM so sánh và chỉ vẽ lại các thành phần thay đổi trên giao diện.

\section{Thiết kế chi tiết}
\label{section:4.2}

\subsection{Thiết kế giao diện}
\label{subsection:4.2.1}

Giao diện hệ thống được thiết kế cho trình duyệt web, hỗ trợ responsive trên nhiều kích thước màn hình từ máy tính (độ phân giải tối thiểu 1280$\times$720) đến thiết bị di động (tối thiểu 375$\times$667). Hệ thống hỗ trợ hai ngôn ngữ (Tiếng Việt và Tiếng Anh) thông qua thư viện \texttt{react-i18next}.

Giao diện được tổ chức thành hai nhóm layout riêng biệt. \textbf{UserLayout} dành cho người dùng, gồm thanh điều hướng ngang (navbar) ở trên cùng với các liên kết đến bốn trang: Dashboard, Lịch sử, Bản đồ và Xếp hạng. \textbf{AdminLayout} dành cho quản trị viên, gồm thanh điều hướng dọc (sidebar) bên trái có thể thu gọn, chứa các mục: Dashboard, Gateways, Sensor Nodes, Cảnh báo, Cấu hình, Người dùng, Xuất dữ liệu và Nhật ký.

Hệ thống sử dụng bộ thành phần giao diện tái sử dụng (component library) tự xây dựng, bao gồm 28 component: Button, Modal, Badge, Input, Select, Table, Pagination, Card, Tooltip và các component chuyên biệt như AQIGauge, MetricCard, HealthAdvisory. Tất cả component tuân thủ thiết kế nhất quán: dark theme với bảng màu chính là xanh dương đậm, các trạng thái tương tác (hover, focus, disabled) được định nghĩa thống nhất, và hiệu ứng chuyển đổi mượt mà (transition 200ms).

%todo: Wireframe/mockup 3–4 màn hình chính — user tự thêm hình

\subsection{Thiết kế lớp}
\label{subsection:4.2.2}

Phần này trình bày thiết kế chi tiết cho bốn module quan trọng nhất trên Backend Server.

\textbf{Module telemetryService} chịu trách nhiệm xử lý dữ liệu đo lường từ Gateway. Module này cung cấp hai hàm chính: \texttt{processTelemetry(gateway\_id, data)} xử lý gói dữ liệu batch và \texttt{processHeartbeat(gateway\_id)} xử lý tín hiệu heartbeat. Hàm \texttt{processTelemetry} thực hiện tuần tự ba bước: lọc trùng lặp qua module \texttt{dedupCache}, thực thi transaction cập nhật trạng thái thiết bị và lưu dữ liệu đo lường, sau đó kiểm tra ngưỡng cảnh báo ngoài transaction.

\textbf{Module aqiService} triển khai thuật toán tính \gls{AQI} theo chuẩn \gls{USEPA} đã trình bày ở Chương 3. Module cung cấp năm hàm: \texttt{calculateSubAQI} tính AQI thành phần cho một chất ô nhiễm, \texttt{calculateAQI} tính AQI tổng hợp (giá trị lớn nhất giữa PM2.5 và PM10), \texttt{getAQIInfo} ánh xạ giá trị AQI sang mức và mã màu, \texttt{getCO2Info} và \texttt{getTVOCInfo} đánh giá mức CO\textsubscript{2} và TVOC. Bảng breakpoint được lưu trữ dưới dạng mảng hằng số trong mã nguồn, tương ứng chính xác với Bảng \ref{table:bp_pm25} và Bảng \ref{table:bp_pm10} ở Chương 3.

\textbf{Module alertService} quản lý toàn bộ vòng đời cảnh báo. Module cung cấp sáu hàm: \texttt{getAlerts} truy vấn danh sách cảnh báo với hỗ trợ lọc và phân trang, \texttt{acknowledgeAlert} xác nhận xử lý, \texttt{deleteAlert} xóa cảnh báo, \texttt{checkThresholdsAndCreateAlerts} kiểm tra năm thông số (PM2.5, PM10, CO\textsubscript{2}, TVOC, nhiệt độ) so với ngưỡng cấu hình và tạo cảnh báo khi vượt ngưỡng, \texttt{createConnectivityAlerts} tạo cảnh báo mất kết nối cho thiết bị offline, và \texttt{cleanupOldAlerts} xóa cảnh báo quá 30 ngày. Cơ chế cooldown 15 phút ngăn chặn tạo cảnh báo trùng lặp cho cùng thiết bị và cùng thông số trong khoảng thời gian ngắn.

\textbf{Module stationService} cung cấp dữ liệu cho các trang người dùng. Hàm \texttt{getNearestStation} sử dụng PostGIS để truy vấn trạm đo gần nhất dựa trên toạ độ GPS của người dùng thông qua hàm \texttt{ST\_Distance}. Hàm \texttt{getStationHistory} truy vấn dữ liệu lịch sử từ Continuous Aggregate \texttt{hourly\_measurements} cho chế độ xem theo giờ, hoặc từ bảng \texttt{measurements} gốc cho chế độ xem theo ngày.

%todo: 2–3 biểu đồ trình tự (sequence diagram) cho use case quan trọng — user tự thêm hình

\subsection{Thiết kế cơ sở dữ liệu}
\label{subsection:4.2.3}

Cơ sở dữ liệu được thiết kế trên PostgreSQL 15 với hai extension: TimescaleDB cho dữ liệu chuỗi thời gian và PostGIS cho dữ liệu không gian. Lược đồ gồm bảy bảng chính và một materialized view.

%todo: E-R diagram cho 7 bảng — user tự thêm hình

\textbf{Bảng users} lưu thông tin tài khoản với khóa chính là UUID tự sinh, bao gồm username (duy nhất), password\_hash (mã hóa bcrypt) và role (admin hoặc user).

\textbf{Bảng gateways} lưu thông tin Gateway với khóa chính là mã định danh dạng chuỗi (ví dụ ``GW\_001''), bao gồm tên, mô tả vị trí, trạng thái online/offline và thời điểm gửi dữ liệu cuối (last\_seen).

\textbf{Bảng sensor\_nodes} lưu thông tin Sensor Node với khóa ngoại tham chiếu đến bảng gateways (quan hệ 1-N: một Gateway quản lý nhiều Sensor Node, xóa Gateway sẽ xóa cascade các Sensor Node). Cột \texttt{geom} kiểu \texttt{geometry(Point, 4326)} lưu toạ độ GPS theo hệ tọa độ WGS84, cho phép truy vấn không gian qua PostGIS. Các cột \texttt{battery\_level} và \texttt{lora\_rssi} lưu thông tin trạng thái phần cứng.

\textbf{Bảng measurements} là hypertable TimescaleDB, được tự động phân mảnh theo cột thời gian \texttt{time}. Khóa chính composite gồm (\texttt{time}, \texttt{node\_id}) đảm bảo mỗi Sensor Node chỉ có một bản ghi tại mỗi thời điểm. Bảng lưu sáu thông số: PM2.5, PM10, CO\textsubscript{2}, TVOC, nhiệt độ và độ ẩm. Retention policy tự động xóa dữ liệu thô quá 3 tháng.

\textbf{Materialized view hourly\_measurements} là Continuous Aggregate tự động tính giá trị trung bình theo giờ cho mỗi Sensor Node. View này được TimescaleDB tự động cập nhật khi có dữ liệu mới ghi vào bảng \texttt{measurements}, đáp ứng yêu cầu phi chức năng truy vấn lịch sử 30 ngày trong vòng 3 giây (Bảng \ref{table:nfr}).

\textbf{Bảng alert\_configs} lưu cấu hình ngưỡng cảnh báo dưới dạng singleton row (chỉ một bản ghi với id = ``default''). Cấu hình bao gồm ngưỡng warn và danger cho bốn thông số (PM2.5, PM10, CO\textsubscript{2}, TVOC), ngưỡng nhiệt độ tối thiểu/tối đa, và chu kỳ lấy mẫu.

\textbf{Bảng alerts} lưu cảnh báo tự động với hai loại: ``threshold'' (vượt ngưỡng) và ``connectivity'' (mất kết nối). Mỗi cảnh báo ghi nhận thông số vi phạm (metric), giá trị đo được (value), ngưỡng so sánh (threshold) và trạng thái xác nhận (acknowledged). Ba chỉ mục trên cột \texttt{node\_id}, \texttt{created\_at} và \texttt{acknowledged} tối ưu hiệu năng truy vấn lọc và phân trang.

\textbf{Bảng audit\_logs} ghi lại mọi hành động quản trị, bao gồm người thực hiện (user\_id, username), loại thao tác (action: CREATE, UPDATE, DELETE), đối tượng thao tác (resource, resource\_id), chi tiết thay đổi (details dạng JSON) và địa chỉ IP. Ba chỉ mục trên cột \texttt{user\_id}, \texttt{created\_at} và \texttt{resource} phục vụ truy vấn tra cứu nhật ký.

\section{Xây dựng ứng dụng}
\label{section:4.3}

\subsection{Thư viện và công cụ sử dụng}
\label{subsection:4.3.1}

Bảng \ref{table:tools_backend} và Bảng \ref{table:tools_frontend} liệt kê các công cụ, thư viện và phiên bản sử dụng để phát triển hệ thống, phân nhóm theo Backend và Frontend.

\begin{table}[H]
\centering
\caption{Danh sách thư viện và công cụ phía Backend}
\label{table:tools_backend}
\small
\begin{tabular}{|p{3.5cm}|p{3.5cm}|p{2cm}|p{3.5cm}|}
\hline
\textbf{Mục đích} & \textbf{Công cụ} & \textbf{Phiên bản} & \textbf{Địa chỉ URL} \\
\hline
Runtime & Node.js & 22.x & https://nodejs.org/ \\
\hline
Web framework & Express & 4.21.0 & https://expressjs.com/ \\
\hline
ORM & Prisma Client & 7.5.0 & https://www.prisma.io/ \\
\hline
Giao tiếp real-time & Socket.IO & 4.7.5 & https://socket.io/ \\
\hline
Xác thực JWT & jsonwebtoken & 9.0.3 & https://github.com/auth0/node-jsonwebtoken \\
\hline
Mã hóa mật khẩu & bcryptjs & 3.0.3 & https://github.com/nicolo-ribaudo/bcrypt.js \\
\hline
Bảo mật HTTP headers & Helmet & 8.1.0 & https://helmetjs.github.io/ \\
\hline
Tác vụ định kỳ & node-cron & 4.2.1 & https://github.com/node-cron/node-cron \\
\hline
Xuất CSV & json2csv & 6.0.0 & https://mircoithink.github.io/json2csv \\
\hline
CSDL & PostgreSQL & 15 & https://www.postgresql.org/ \\
\hline
Extension time-series & TimescaleDB & - & https://www.timescale.com/ \\
\hline
Extension geospatial & PostGIS & - & https://postgis.net/ \\
\hline
\end{tabular}
\end{table}

\begin{table}[H]
\centering
\caption{Danh sách thư viện và công cụ phía Frontend}
\label{table:tools_frontend}
\small
\begin{tabular}{|p{3.5cm}|p{3.5cm}|p{2cm}|p{3.5cm}|}
\hline
\textbf{Mục đích} & \textbf{Công cụ} & \textbf{Phiên bản} & \textbf{Địa chỉ URL} \\
\hline
UI framework & React & 19.2.4 & https://react.dev/ \\
\hline
Build tool & Vite & 8.0.4 & https://vite.dev/ \\
\hline
Routing & React Router & 7.14.0 & https://reactrouter.com/ \\
\hline
HTTP client & Axios & 1.14.0 & https://axios-http.com/ \\
\hline
Bản đồ tương tác & Leaflet + React-Leaflet & 1.9.4 / 5.0.0 & https://react-leaflet.js.org/ \\
\hline
Biểu đồ & ECharts & 6.0.0 & https://echarts.apache.org/ \\
\hline
Đa ngôn ngữ & react-i18next & 17.0.2 & https://react.i18next.com/ \\
\hline
Quản lý trạng thái & Zustand & 5.0.12 & https://zustand.docs.pmnd.rs/ \\
\hline
Icon & Lucide React & 1.7.0 & https://lucide.dev/ \\
\hline
Socket client & socket.io-client & 4.8.3 & https://socket.io/ \\
\hline
\end{tabular}
\end{table}

Ngoài ra, hệ thống sử dụng các công cụ phát triển và triển khai: VS Code làm IDE chính, PlatformIO với framework Arduino để phát triển firmware ESP32, Docker và Docker Compose để container hóa, Nginx làm reverse proxy, và Certbot để quản lý chứng chỉ SSL.

\subsection{Kết quả đạt được}
\label{subsection:4.3.2}

Sản phẩm cuối cùng bao gồm bốn thành phần được đóng gói riêng biệt. Thứ nhất, \textbf{Firmware Sensor Node} chạy trên ESP32, sử dụng FreeRTOS để quản lý đồng thời các tác vụ đọc cảm biến, gửi LoRa và hiển thị OLED, với ba biến thể: board 38 chân, board 30 chân, và board 30 chân có chế độ deep-sleep. Thứ hai, \textbf{Firmware Gateway} chạy trên ESP32, sử dụng kiến trúc superloop để nhận gói LoRa và chuyển tiếp lên Backend Server qua WiFi, với hai biến thể cho board 38 chân và 30 chân. Thứ ba, \textbf{Backend Server} là ứng dụng Node.js/Express cung cấp \gls{REST} \gls{API} và giao tiếp real-time qua Socket.IO. Thứ tư, \textbf{Frontend} là ứng dụng React \gls{SPA} được build bởi Vite thành các tệp tĩnh, phục vụ bởi Nginx.

Bảng \ref{table:stats} thống kê quy mô mã nguồn của từng thành phần.

\begin{table}[H]
\centering
\caption{Thống kê quy mô mã nguồn hệ thống}
\label{table:stats}
\small
\begin{tabular}{|p{4cm}|p{3cm}|p{3cm}|p{2.5cm}|}
\hline
\textbf{Thành phần} & \textbf{Số file mã nguồn} & \textbf{Ngôn ngữ chính} & \textbf{Ghi chú} \\
\hline
Firmware Sensor Node & 3 biến thể & C/C++ (Arduino) & FreeRTOS \\
\hline
Firmware Gateway & 2 biến thể & C/C++ (Arduino) & Superloop \\
\hline
Backend Server & 14 controllers, 15 services & JavaScript (Node.js) & Express 4.21 \\
\hline
Frontend & 16 pages, 28 components & JavaScript (React) & Vite 8.0 \\
\hline
Database & 7 bảng, 1 view & SQL & TimescaleDB + PostGIS \\
\hline
Triển khai & 4 containers & YAML, Nginx conf & Docker Compose \\
\hline
\end{tabular}
\end{table}

\subsection{Minh họa các chức năng chính}
\label{subsection:4.3.3}

Phần này minh hoạ các chức năng chính của hệ thống thông qua giao diện thực tế trên môi trường production.

\textbf{Dashboard chất lượng không khí.} Khi người dùng truy cập trang Dashboard, hệ thống yêu cầu quyền truy cập vị trí, sau đó hiển thị dữ liệu của trạm đo gần nhất. Giao diện hiển thị chỉ số AQI tổng hợp dưới dạng gauge lớn ở trung tâm, kèm mã màu theo chuẩn \gls{USEPA}. Phía dưới là sáu thẻ metric hiển thị PM2.5, PM10, CO\textsubscript{2}, TVOC, nhiệt độ và độ ẩm, mỗi thẻ kèm đánh giá mức độ. Bên cạnh là khuyến nghị sức khỏe dựa trên mức AQI hiện tại. Dữ liệu được cập nhật thời gian thực khi Backend phát sự kiện qua Socket.IO.

%todo: Screenshot Dashboard — user tự thêm hình

\textbf{Lịch sử dữ liệu.} Trang lịch sử cho phép người dùng chọn trạm đo, thông số cần xem và chế độ hiển thị (theo giờ hoặc theo ngày). Biểu đồ được vẽ bằng ECharts, hỗ trợ zoom, tooltip và responsive. Chế độ theo giờ truy vấn từ Continuous Aggregate \texttt{hourly\_measurements} để đảm bảo hiệu năng.

%todo: Screenshot Lịch sử dữ liệu — user tự thêm hình

\textbf{Bản đồ AQI.} Trang bản đồ hiển thị toàn bộ các trạm đo trên bản đồ Leaflet với tile OpenStreetMap. Mỗi trạm được đánh dấu bằng marker hình tròn với mã màu AQI. Người dùng nhấn vào marker để xem popup thông tin chi tiết gồm tên trạm, giá trị AQI, PM2.5 và thời gian cập nhật cuối.

%todo: Screenshot Bản đồ AQI — user tự thêm hình

\textbf{Xếp hạng ô nhiễm.} Trang xếp hạng hiển thị danh sách tất cả trạm đo được sắp xếp theo mức AQI từ cao đến thấp. Mỗi dòng hiển thị tên trạm, giá trị AQI, nồng độ PM2.5 và badge mã màu tương ứng.

%todo: Screenshot Xếp hạng — user tự thêm hình

\textbf{Quản lý thiết bị.} Trang quản trị thiết bị hiển thị hai bảng riêng biệt cho Gateway và Sensor Node. Mỗi bảng cho biết trạng thái online/offline, thời gian gửi cuối, mức pin (đối với Sensor Node) và cường độ tín hiệu LoRa (RSSI). Quản trị viên thêm, sửa, xóa thiết bị thông qua modal form với xác thực dữ liệu đầu vào.

%todo: Screenshot Quản lý thiết bị — user tự thêm hình

\textbf{Giám sát cảnh báo.} Trang cảnh báo hiển thị danh sách phân trang các cảnh báo, hỗ trợ lọc theo thiết bị, mức độ (warn/danger) và trạng thái xác nhận. Mỗi cảnh báo hiển thị thông tin thời gian, thiết bị, thông số, giá trị vi phạm và ngưỡng. Quản trị viên có thể xác nhận đã xử lý (acknowledge) từng cảnh báo. Cảnh báo mới được phát thời gian thực qua Socket.IO.

%todo: Screenshot Giám sát cảnh báo — user tự thêm hình

\textbf{Cấu hình hệ thống.} Trang cấu hình cho phép quản trị viên thiết lập ngưỡng cảnh báo cho năm thông số, mỗi thông số có hai mức warn và danger. Giao diện hiển thị các trường nhập liệu nhóm theo thông số, kèm giá trị mặc định theo chuẩn \gls{USEPA}.

%todo: Screenshot Cấu hình — user tự thêm hình

\textbf{Xuất dữ liệu CSV.} Trang xuất dữ liệu cho phép quản trị viên chọn Sensor Node và khoảng thời gian, sau đó xuất tệp CSV chứa toàn bộ dữ liệu đo lường. Tệp CSV được tạo phía Backend bằng thư viện \texttt{json2csv} và trả về dưới dạng download.

%todo: Screenshot Xuất CSV — user tự thêm hình

\section{Kiểm thử}
\label{section:4.4}

Hệ thống được kiểm thử bằng phương pháp kiểm thử chức năng thủ công (manual functional testing), thiết kế các trường hợp kiểm thử dựa trên đặc tả use case ở Chương 2. Phần này trình bày kết quả kiểm thử cho ba chức năng quan trọng nhất.

\subsection{Kiểm thử chức năng Xem dashboard}
\label{subsection:4.4.1}

Bảng \ref{table:test_dashboard} trình bày các trường hợp kiểm thử cho chức năng xem dashboard chất lượng không khí.

\begin{table}[H]
\centering
\caption{Trường hợp kiểm thử chức năng Xem dashboard}
\label{table:test_dashboard}
\small
\begin{tabular}{|p{0.5cm}|p{3.5cm}|p{4cm}|p{2.5cm}|p{1.5cm}|}
\hline
\textbf{STT} & \textbf{Mô tả} & \textbf{Kết quả mong đợi} & \textbf{Kết quả thực tế} & \textbf{Đạt} \\
\hline
1 & Truy cập dashboard, cho phép truy cập vị trí & Hiển thị dữ liệu trạm gần nhất & Đúng & Đạt \\
\hline
2 & Truy cập dashboard, từ chối truy cập vị trí & Hiển thị dữ liệu trạm mặc định & Đúng & Đạt \\
\hline
3 & Dashboard nhận dữ liệu mới qua Socket.IO & Các thẻ metric cập nhật tự động & Đúng & Đạt \\
\hline
4 & Không có trạm nào hoạt động & Hiển thị thông báo ``Không có dữ liệu'' & Đúng & Đạt \\
\hline
5 & Kiểm tra mã màu AQI & Gauge và badge hiển thị đúng mã màu theo mức AQI & Đúng & Đạt \\
\hline
\end{tabular}
\end{table}

\subsection{Kiểm thử chức năng Quản lý thiết bị}
\label{subsection:4.4.2}

Bảng \ref{table:test_device} trình bày các trường hợp kiểm thử cho chức năng quản lý thiết bị.

\begin{table}[H]
\centering
\caption{Trường hợp kiểm thử chức năng Quản lý thiết bị}
\label{table:test_device}
\small
\begin{tabular}{|p{0.5cm}|p{3.5cm}|p{4cm}|p{2.5cm}|p{1.5cm}|}
\hline
\textbf{STT} & \textbf{Mô tả} & \textbf{Kết quả mong đợi} & \textbf{Kết quả thực tế} & \textbf{Đạt} \\
\hline
1 & Thêm mới Gateway với thông tin hợp lệ & Gateway được tạo, hiển thị trong danh sách & Đúng & Đạt \\
\hline
2 & Thêm mới Sensor Node, chọn Gateway cha & Sensor Node được tạo, liên kết đúng Gateway & Đúng & Đạt \\
\hline
3 & Sửa thông tin thiết bị & Thông tin cập nhật thành công & Đúng & Đạt \\
\hline
4 & Xóa Gateway có Sensor Node con & Xóa cascade cả Sensor Node con & Đúng & Đạt \\
\hline
5 & Thêm thiết bị với dữ liệu không hợp lệ (để trống tên) & Hiển thị thông báo lỗi & Đúng & Đạt \\
\hline
6 & Kiểm tra audit log sau thao tác & Nhật ký ghi nhận đúng thao tác, người thực hiện, IP & Đúng & Đạt \\
\hline
\end{tabular}
\end{table}

\subsection{Kiểm thử chức năng Giám sát cảnh báo}
\label{subsection:4.4.3}

Bảng \ref{table:test_alert} trình bày các trường hợp kiểm thử cho chức năng giám sát và xử lý cảnh báo.

\begin{table}[H]
\centering
\caption{Trường hợp kiểm thử chức năng Giám sát cảnh báo}
\label{table:test_alert}
\small
\begin{tabular}{|p{0.5cm}|p{3.5cm}|p{4cm}|p{2.5cm}|p{1.5cm}|}
\hline
\textbf{STT} & \textbf{Mô tả} & \textbf{Kết quả mong đợi} & \textbf{Kết quả thực tế} & \textbf{Đạt} \\
\hline
1 & Gửi dữ liệu PM2.5 vượt ngưỡng warn & Cảnh báo mức warn được tạo tự động & Đúng & Đạt \\
\hline
2 & Gửi dữ liệu PM2.5 vượt ngưỡng danger & Cảnh báo mức danger được tạo, phát qua Socket.IO & Đúng & Đạt \\
\hline
3 & Gửi dữ liệu vượt ngưỡng lần 2 trong 15 phút & Không tạo cảnh báo trùng (cooldown) & Đúng & Đạt \\
\hline
4 & Thiết bị không gửi dữ liệu quá thời gian quy định & Cron job tạo cảnh báo mất kết nối & Đúng & Đạt \\
\hline
5 & Lọc cảnh báo theo mức độ & Danh sách chỉ hiển thị cảnh báo đúng mức độ & Đúng & Đạt \\
\hline
6 & Xác nhận cảnh báo (acknowledge) & Trạng thái cập nhật thành ``đã xác nhận'' & Đúng & Đạt \\
\hline
\end{tabular}
\end{table}

\subsection{Tổng kết kiểm thử}
\label{subsection:4.4.4}

Tổng cộng 17 trường hợp kiểm thử đã được thực hiện trên ba chức năng chính: Xem dashboard (5 test case), Quản lý thiết bị (6 test case) và Giám sát cảnh báo (6 test case). Tất cả 17/17 trường hợp đều đạt, tỉ lệ 100\%. Kết quả cho thấy các chức năng chính của hệ thống hoạt động đúng theo đặc tả use case đã định nghĩa ở Chương 2.

\section{Triển khai}
\label{section:4.5}

Hệ thống được triển khai trên máy chủ ảo (VPS) DigitalOcean với cấu hình 1 vCPU, 1 GB RAM, 25 GB SSD, chạy hệ điều hành Ubuntu 22.04 LTS. Domain truy cập là \texttt{datn.thamnguyen.dev}, được cấu hình HTTPS thông qua chứng chỉ SSL miễn phí từ Let's Encrypt.

%todo: Biểu đồ triển khai (deployment diagram) — user tự thêm hình

Mô hình triển khai sử dụng Docker Compose để quản lý bốn container. Container \textbf{db} chạy image \texttt{timescale/timescaledb-ha:pg15}, lưu trữ dữ liệu trên volume \texttt{pg\_data} và chỉ mở cổng 5432 trên localhost (truy cập qua SSH tunnel). Container \textbf{backend} chạy ứng dụng Node.js, kết nối đến container db thông qua mạng nội bộ Docker. Container \textbf{nginx} chạy Nginx phục vụ hai chức năng: (i) reverse proxy chuyển tiếp các yêu cầu \gls{HTTP} và WebSocket đến Backend, và (ii) phục vụ các tệp tĩnh Frontend đã được build bởi Vite. Container \textbf{certbot} tự động gia hạn chứng chỉ SSL mỗi 12 giờ.

Quy trình triển khai được thực hiện bằng một lệnh duy nhất: \texttt{docker compose --env-file .env.production -f docker-compose.prod.yml up --build -d}. Docker Compose tự động build image cho Backend và Nginx từ Dockerfile, khởi tạo cơ sở dữ liệu từ \texttt{init.sql}, và kết nối các container qua mạng nội bộ \texttt{app-network}. Toàn bộ cấu hình nhạy cảm (mật khẩu database, JWT secret) được lưu trong file \texttt{.env.production} không đưa lên kho mã nguồn.

%todo: Screenshot/thông tin triển khai thực tế — user tự thêm

Chương này đã trình bày toàn bộ quá trình thiết kế, xây dựng, kiểm thử và triển khai hệ thống giám sát chất lượng không khí. Về thiết kế, hệ thống áp dụng kiến trúc bốn tầng (cảm biến, trung chuyển, máy chủ, hiển thị) kết hợp mô hình Hub-Spoke, tổ chức mã nguồn theo nguyên tắc phân lớp rõ ràng. Cơ sở dữ liệu tận dụng TimescaleDB cho dữ liệu chuỗi thời gian và PostGIS cho truy vấn không gian. Về xây dựng, hệ thống gồm bốn thành phần (firmware, backend, frontend, database) với tổng cộng 16 trang giao diện và 28 component tái sử dụng. Về kiểm thử, 17/17 trường hợp kiểm thử chức năng đều đạt. Về triển khai, hệ thống chạy ổn định trên VPS DigitalOcean với Docker Compose. Trên cơ sở kết quả đạt được, Chương 5 sẽ phân tích các giải pháp và đóng góp nổi bật của đồ án.

\end{document}
